\subsection*{1.4. 一个特殊情形}

本讨论班后续不会用到本节。它将第 XIII 讲引言中概述的一个构造形式化。

\subsubsection*{1.4.1.}
设 $D$ 是有定向圆盘，$0\in D$，并令 $D^*=D-\{0\}$。$\mathcal D^*$ 上的一个\underline{拓扑空间} $\mathcal X^*$ 是一个函子
\[
(\text{万有覆盖，底空间为 }D^*)
\longrightarrow(\text{拓扑空间}).
\]
选择一个万有覆盖 $\widetilde D^*$ 后，它就是一个带有 $\pi_1(D^*)$ 作用的拓扑空间 $X_\eta$；当 $\widetilde D^*$ 是 $(D^*,t)$ 的万有覆盖时，将其记为 $\mathcal X_t^*$。

$\mathcal X^*$ 上的一个\underline{层} 是一个函子
\[
(\text{万有覆盖，底空间为 }D^*)
\longrightarrow
(\text{带有层的拓扑空间}),
\]
或等价地，是 $X_\eta$ 上的一个 $\pi_1(D^*)$-等变层。

对 $\mathcal X^*$，赋予 $D^*$ 上的拓扑空间
\[
\beta^*\mathcal X^*
=X_\eta\times\widetilde D^*/\pi_1(D^*)
\]
$\mathcal X^*$ 上的每个层 $F$ 下降为 $\beta^*\mathcal X^*$ 上的层 $\beta^*F$，给出拓扑斯的交换图
\[
\begin{tikzcd}
\beta^*\mathcal X^* \arrow[r] \arrow[d]
  & \mathcal X^* \arrow[d] \\
D^* \arrow[r, "\beta"]
  & \mathcal D^*.
\end{tikzcd}
\]

\subsubsection*{1.4.2.}
$\mathcal D$ 上的一个\underline{拓扑空间} $\mathcal X$ 由以下数据组成：
\begin{enumerate}[label=(\alph*),leftmargin=2.7em]
\item $\mathcal D^*$ 上的拓扑空间 $\mathcal X^*$；
\item 拓扑空间 $X_0$；
\item $\pi_1(D^*)$-不变态射
\[
\mathrm{sp}:X_\eta\longrightarrow X_0,
\]
其中 $\pi_1(D^*)$ 在 $X_0$ 上平凡作用。
\end{enumerate}

$\mathcal X$ 上的一个\underline{层} 由 $X_0$ 上的层 $F_0$、$\mathcal X^*$ 上由 $F_\eta$ 给出的层，以及一个等变态射
\[
\mathrm{sp}:\mathrm{sp}^*F_0\longrightarrow F_\eta
\qquad
(\text{或 }F_0\longrightarrow\mathrm{sp}_*F_\eta).
\]
组成。有交换图
\[
\begin{tikzcd}
X_0 \arrow[r, hook] \arrow[d]
  & \mathcal X \arrow[d]
  & \mathcal X^* \arrow[l, hook] \arrow[d] \\
\{0\} \arrow[r, hook]
  & \mathcal D
  & \mathcal D^* \arrow[l, hook]
\end{tikzcd}
\]
此外，还有一个拓扑斯态射
\[
\Psi:\mathcal X\longrightarrow X_0\times\mathcal D.
\]
它由前述构造的函子性应用于
\[
(X_\eta,X_0,\mathrm{sp})
\longrightarrow(X_0,X_0,\mathrm{Id}),
\]
而得，且满足
\[
\Psi(F_0,F_\eta,\mathrm{sp})
=(F_0,\mathrm{sp}_*F_\eta,\mathrm{sp}).
\]
此态射诱导
\[
\Psi_\eta:\mathcal X^*\longrightarrow X_0\times\mathcal D^*.
\]

